\documentclass[10pt,a4paper]{amsart}
\usepackage[margin=19mm]{geometry}
\usepackage{amsmath,amssymb,mathtools}
\usepackage[scr=rsfs,cal=euler]{mathalfa}
\usepackage{xfrac}
\usepackage[unicode,hidelinks,bookmarks=false]{hyperref}
\newcommand{\ie}{i.e.\ }
\newcommand{\cf}{cf.\ }
\newcommand{\resp}{respectively\ }
\newcommand{\Subsection}{\subsection}
\providecommand{\nobreakdash}{\nobreak\hbox{-}}
\setlength{\emergencystretch}{2em}
\multlinegap0pt
\usepackage{booktabs}
\usepackage{algorithm}\usepackage{algpseudocode}
\newtheorem{theorem}{Theorem}
\renewcommand{\d}{\ensuremath{\,\mathrm{d}}}
\newcommand{\essinf}{\operatorname{ess\,inf}}         % Essential supremum
\newcommand{\qp}[1]{\ensuremath{\!\left({#1}\right)}}
\renewcommand{\vec}[1]{\ensuremath{\boldsymbol{#1}}}
\title{Characteristic Sweeps and Source Iteration for Charged-Particle Transport with Continuous Slowing-Down and Angular Scattering}
\author{Ben S. Ashby, Alex Lukyanov, Tristan Pryer}
\date{Source: arXiv:2602.07857v1 (CC BY 4.0)}
\begin{document}
\maketitle

\noindent\emph{Source and licence.} Characteristic Sweeps and Source Iteration for Charged-Particle Transport with Continuous Slowing-Down and Angular Scattering. Version arXiv:2602.07857v1 (CC BY 4.0), distributed by arXiv under the Creative Commons Attribution 4.0 International licence, http://creativecommons.org/licenses/by/4.0/, as recorded in the arXiv licence metadata for this exact version and shown on the abstract page https://arxiv.org/abs/2602.07857v1. Full licence notice: \texttt{licenses/arxiv-2602.07857-CC-BY-4.0.txt}.\par\smallskip

\noindent\textit{Editorial scope.} Counted unit: the article's theorem the:coercivity (source0.tex lines 871--934). The article's complete original proof of it is delivered with it (source0.tex lines 936--1074, one \texttt{proof} environment of 139 lines). No mathematics is written, repaired or re-derived: the statement and the proof are the article's own lines, in the article's own notation, and no retained line is substituted or re-flowed.  Retained uncounted as the source-local context the proof needs: lines 775--805.  Nothing was omitted from any retained span, so the package carries no omission record.  No retained line contains a citation, so the wrapper carries no bibliography and no bibliography entry.  No retained line contains a figure environment, an image-inclusion command or drawing code, so no image is delivered and the package declares no asset dependency.  Editorial formatting only, in addition: an AMS \texttt{amsart} preamble that restates the article's own declarations and macros used by the retained lines, namely the environments \texttt{theorem} on separate counters, together with the standard packages those lines need.  The retained environments print in this document as Theorem 1 in order of appearance, whereas the article prints the same items with the numbering of its own full document.  The complete native source of the article as distributed in the arXiv e-print for this exact version is bundled unchanged as \texttt{sources/194218/source0.tex}.
\begin{equation}\label{eq:inflowBC}
  \psi = g \quad \text{on }\Gamma_-.
\end{equation}
We assume $g$ is square-integrable with respect to the natural flux
weights induced by the boundary term $\vec\omega\cdot\vec n_{\vec x}-S(E)n_E$,
namely $|\vec\omega\cdot\vec n_{\vec x}|$ on $\Gamma_-^{\vec x}$ and
$S(E_{\max})$ on $\Gamma_-^{E}$:
\[
  g\in L^2(\Gamma_-^{\vec x}, |\vec\omega\cdot \vec n_{\vec x}| \d\gamma_{\vec x} \d E \d\vec\omega)
  \ \cap\
  L^2(\Gamma_-^{E}, S(E_{\max}) \d\vec x \d\vec\omega).
\]
No boundary condition is imposed on $\Gamma_+$.

\section{Variational Formulation and Functional Setting}
\label{sec:variational}

To analyse well-posedness and enable discretisation, we adopt a
variational formulation. The first-order transport and energy-advection
terms, together with the non-symmetric gain operator, motivate a graph
space that captures directional and energy drift.

We work in the graph space
\begin{equation}
  \label{eq:Uspace}
  \mathcal{U}
  :=
  \Big\{ u \in L^2(\Omega)  : 
  \vec{\omega}\cdot\nabla_x u \in L^2(\Omega),
  \ \partial_E\qp{S(E)u} \in L^2(\Omega) \Big\}.
\end{equation}

\begin{theorem}[Coercivity and boundedness of the continuous bilinear forms]
\label{the:coercivity}
Assume $S\in W^{1,\infty}(I)$ with $S(E)>0$ and $S'(E)\le 0$ on $I$,
$\sigma_T\in L^\infty(I)$ with $\sigma_T(E)\ge 0$, and $\sigma_S\in
L^\infty(\mathbb{S}^{d-1}\times\mathbb{S}^{d-1}\times I\times I)$ with
$\sigma_S\ge 0$.

Define the coercivity norm
\begin{equation*}
  \|u\|_{\mathcal{U}}^2
  :=
  \|\sigma_T^{1/2}u\|_{L^2(\Omega)}^2
  + \|(-S'(E))^{1/2}u\|_{L^2(\Omega)}^2
  + \int_{\Gamma_+} \qp{\vec{\omega}\cdot\vec{n}_{\vec x} - S(E) n_E} u^2 \d s,
\end{equation*}
and the stronger norm
\begin{equation*}
  \|u\|_{\mathcal{U}^*}^2
  :=
  \|u\|_{\mathcal{U}}^2
  + \|\nabla_{\vec x} u\|_{L^2(\Omega)}^2
  + \|\partial_E u\|_{L^2(\Omega)}^2.
\end{equation*}
Assume moreover that
\begin{equation*}
  \mu := \essinf_{E\in I}\qp{-S'(E)} > 0,
\end{equation*}
so that $\|w\|_{L^2(\Omega)} \le
\mu^{-1/2}\|(-S')^{1/2}w\|_{L^2(\Omega)} \le
\mu^{-1/2}\|w\|_{\mathcal{U}}$.

Then the bilinear forms $a(u,v;\vec{\omega})$ and $s(u,v;\vec{\omega})$
defined in Section~\ref{sec:variational} satisfy the following:

\begin{enumerate}
\item
  Coercivity. If $u=0$ on $\Gamma_-$, then
  \begin{equation*}
    \int_{\mathbb{S}^{d-1}} a(u,u;\vec{\omega}) \d\vec{\omega}
    \ge
    \tfrac{1}{2}\|u\|_{\mathcal{U}}^2.
  \end{equation*}
  
\item Boundedness. There exists $C_a>0$, depending only on
  $\|S\|_{W^{1,\infty}(I)}$, $\|\sigma_T\|_{L^\infty(I)}$ and $\mu$, such
  that for all $u\in\mathcal{U}^*$ and all $v\in\mathcal{U}$,
  \begin{equation*}
    \Big|\int_{\mathbb{S}^{d-1}} a(u,v;\vec{\omega}) \d\vec{\omega}\Big|
    \le
    C_a \|u\|_{\mathcal{U}^*}\|v\|_{\mathcal{U}}.
  \end{equation*}
  Moreover there exists $C_s>0$, depending only on
  $\|\sigma_S\|_{L^\infty}$ and the measures of $I$ and $\mathbb{S}^{d-1}$,
  such that
  \begin{equation*}
    \Big|\int_{\mathbb{S}^{d-1}} s(u,v;\vec{\omega}) \d\vec{\omega}\Big|
    \le
    C_s \|u\|_{L^2(\Omega)}\|v\|_{L^2(\Omega)}
    \le
    C_s \mu^{-1}\|u\|_{\mathcal{U}}\|v\|_{\mathcal{U}}.
  \end{equation*}
\end{enumerate}

\end{theorem}

\begin{proof}
Fix $\vec{\omega}\in\mathbb{S}^{d-1}$ and abbreviate
$a(u,u):=a(u,u;\vec{\omega})$. By definition,
\begin{align*}
  a(u,u)
  &=
  \int_D \int_I
  \Big[\qp{\vec{\omega}\cdot\nabla_{\vec x} u} u
  - \partial_E\qp{S u} u
  + \sigma_T u^2\Big]\d E \d\vec{x}
  + \tfrac{1}{2}\int_{\Gamma_+(\vec\omega)}\qp{\vec{\omega}\cdot\vec{n}_{\vec x} - S n_E} u^2 \d\gamma.
\end{align*}

We first integrate by parts in $\vec{x}$. Since
$\qp{\vec{\omega}\cdot\nabla_{\vec x} u} u = \tfrac{1}{2}\vec{\omega}\cdot\nabla_{\vec x}\qp{u^2}$,
\begin{equation*}
  \int_D \int_I \qp{\vec{\omega}\cdot\nabla_{\vec x} u} u \d E \d\vec{x}
  =
  \tfrac{1}{2}\int_{\partial D\times I}\qp{\vec{\omega}\cdot\vec{n}_{\vec x}} u^2 \d\gamma.
\end{equation*}
For the energy-advection term, write
\begin{equation*}
  -\partial_E\qp{S u} u
  =
  -\tfrac{1}{2}\partial_E\qp{S u^2} + \tfrac{1}{2}\qp{-S'} u^2,
\end{equation*}
so that integration by parts in $E$ gives
\begin{equation*}
  -\int_D \int_I \partial_E\qp{S u} u \d E \d\vec{x}
  =
  -\tfrac{1}{2}\int_{D\times\partial I} S n_E u^2 \d\gamma
  + \tfrac{1}{2}\int_D \int_I \qp{-S'(E)} u^2 \d E \d\vec{x}.
\end{equation*}
Combining the $\vec{x}$ and $E$ boundary contributions yields
\begin{equation*}
  \tfrac{1}{2}\int_{\partial(D\times I)}
  \qp{\vec{\omega}\cdot\vec{n}_{\vec x} - S n_E} u^2 \d\gamma
  =
  \tfrac{1}{2}\int_{\Gamma_+(\vec\omega)}\qp{\vec{\omega}\cdot\vec{n}_{\vec x} - S n_E} u^2 \d\gamma
  + \tfrac{1}{2}\int_{\Gamma_-(\vec\omega)}\qp{\vec{\omega}\cdot\vec{n}_{\vec x} - S n_E} u^2 \d\gamma.
\end{equation*}
Therefore
\begin{align*}
  a(u,u)
  &=
  \|\sigma_T^{1/2}u(\cdot,\cdot,\vec\omega)\|_{L^2(D\times I)}^2
  + \tfrac{1}{2}\|(-S')^{1/2}u(\cdot,\cdot,\vec\omega)\|_{L^2(D\times I)}^2
  \\
  &\qquad + \int_{\Gamma_+(\vec\omega)}\qp{\vec{\omega}\cdot\vec{n}_{\vec x} - S n_E} u^2 \d\gamma
  + \tfrac{1}{2}\int_{\Gamma_-(\vec\omega)}\qp{\vec{\omega}\cdot\vec{n}_{\vec x} - S n_E} u^2 \d\gamma.
\end{align*}
If $u=0$ on $\Gamma_-$, then for each fixed $\vec\omega$ we have $u=0$
on $\Gamma_-(\vec\omega)$ and the last term vanishes. Hence
\begin{equation*}
  a(u,u)
  \ge
  \tfrac{1}{2}\|\sigma_T^{1/2}u(\cdot,\cdot,\vec\omega)\|_{L^2(D\times I)}^2
  + \tfrac{1}{2}\|(-S')^{1/2}u(\cdot,\cdot,\vec\omega)\|_{L^2(D\times I)}^2
  + \tfrac{1}{2}\int_{\Gamma_+(\vec\omega)}\qp{\vec{\omega}\cdot\vec{n}_{\vec x} - S n_E} u^2 \d\gamma.
\end{equation*}
Integrating this inequality over $\mathbb{S}^{d-1}$ and using $\d
s=\d\gamma \d\vec\omega$ gives
\begin{equation*}
  \int_{\mathbb{S}^{d-1}} a(u,u;\vec{\omega}) \d\vec{\omega}
  \ge
  \tfrac{1}{2}\|\sigma_T^{1/2}u\|_{L^2(\Omega)}^2
  + \tfrac{1}{2}\|(-S')^{1/2}u\|_{L^2(\Omega)}^2
  + \tfrac{1}{2}\int_{\Gamma_+}\qp{\vec{\omega}\cdot\vec{n}_{\vec x} - S n_E} u^2 \d s
  =
  \tfrac{1}{2}\|u\|_{\mathcal{U}}^2,
\end{equation*}
which proves coercivity.

For boundedness of $a$, fix $\vec{\omega}$ and apply Cauchy--Schwarz,
together with
\[
\|\partial_E\qp{S u}\|_{L^2(\Omega)}\le \|S\|_{L^\infty(I)}\|\partial_E u\|_{L^2(\Omega)} + \|S'\|_{L^\infty(I)}\|u\|_{L^2(\Omega)}
\]
we have
\begin{align*}
  \Big|\int_D\int_I \qp{\vec{\omega}\cdot\nabla_{\vec x} u} v\Big|
  &\le
  \|\nabla_{\vec x} u(\cdot,\cdot,\vec\omega)\|_{L^2(D\times I)}\|v(\cdot,\cdot,\vec\omega)\|_{L^2(D\times I)},\\
  \Big|\int_D\int_I \partial_E\qp{S u} v\Big|
  &\le
  \|\partial_E\qp{S u}(\cdot,\cdot,\vec\omega)\|_{L^2(D\times I)}\|v(\cdot,\cdot,\vec\omega)\|_{L^2(D\times I)}\\
  &\le
  \|S\|_{W^{1,\infty}(I)}\qp{\|\partial_E u(\cdot,\cdot,\vec\omega)\|_{L^2(D\times I)}+\|u(\cdot,\cdot,\vec\omega)\|_{L^2(D\times I)}}\|v(\cdot,\cdot,\vec\omega)\|_{L^2(D\times I)},\\
  \Big|\int_D\int_I \sigma_T u v\Big|
  &\le
  \|\sigma_T\|_{L^\infty(I)}\|u(\cdot,\cdot,\vec\omega)\|_{L^2(D\times I)}\|v(\cdot,\cdot,\vec\omega)\|_{L^2(D\times I)},\\
  \Big|\tfrac{1}{2}\int_{\Gamma_+(\vec\omega)}\qp{\vec{\omega}\cdot\vec{n}_{\vec x} - S n_E} u v \d\gamma\Big|
  &\le
  \tfrac{1}{2}
  \Big(\int_{\Gamma_+(\vec\omega)}\qp{\vec{\omega}\cdot\vec{n}_{\vec x} - S n_E} u^2 \d\gamma\Big)^{1/2}
  \Big(\int_{\Gamma_+(\vec\omega)}\qp{\vec{\omega}\cdot\vec{n}_{\vec x} - S n_E} v^2 \d\gamma\Big)^{1/2}.
\end{align*}
Using $\|w\|_{L^2(\Omega)}\le \mu^{-1/2}\|w\|_{\mathcal{U}}$ for $w=v$ and
$\|u\|_{L^2(\Omega)}\le \mu^{-1/2}\|u\|_{\mathcal{U}}\le \mu^{-1/2}\|u\|_{\mathcal{U}^*}$,
we obtain a bound of the form
$|a(u,v;\vec{\omega})|\le C_a\|u\|_{\mathcal{U}^*}\|v\|_{\mathcal{U}}$ with $C_a$ independent of $\vec{\omega}$.
Integrating over $\mathbb{S}^{d-1}$ yields the stated estimate.

For the scattering form, define
\begin{equation*}
  \mathcal{K}[u]\qp{\vec{x},E,\vec{\omega}}
  :=
  \int_{\mathbb{S}^{d-1}}\int_I
  \sigma_S\qp{\vec{\omega},\vec{\omega}',E'\to E} 
  u\qp{\vec{x},E',\vec{\omega}'} \d E' \d\vec{\omega}'.
\end{equation*}
By Fubini and $\|\sigma_S\|_{L^\infty}<\infty$,
\begin{equation*}
  |\mathcal{K}[u]\qp{\vec{x},E,\vec{\omega}}|
  \le
  \|\sigma_S\|_{L^\infty}\int_{\mathbb{S}^{d-1}}\int_I
  |u\qp{\vec{x},E',\vec{\omega}'}| \d E' \d\vec{\omega}'.
\end{equation*}
Applying Cauchy--Schwarz in $(E',\vec{\omega}')$ gives
\[
  |\mathcal{K}[u]\qp{\vec{x},E,\vec{\omega}}|
  \le
  \|\sigma_S\|_{L^\infty} \qp{|I| |\mathbb S^{d-1}|}^{1/2}
  \Big(\int_{\mathbb{S}^{d-1}}\int_I |u\qp{\vec{x},E',\vec{\omega}'}|^2 \d E' \d\vec{\omega}'\Big)^{1/2}.
\]
Squaring, integrating over $(\vec{x},E,\vec{\omega})$ and using Fubini yields
$\|\mathcal{K}[u]\|_{L^2(\Omega)} \le C_s\|u\|_{L^2(\Omega)}$ with $C_s$ depending only on
$\|\sigma_S\|_{L^\infty}$ and the measures of $I$ and $\mathbb{S}^{d-1}$. Hence
\begin{equation*}
  \Big|\int_{\mathbb{S}^{d-1}} s(u,v;\vec{\omega}) \d\vec{\omega}\Big|
  =
  \Big|\int_\Omega v \mathcal{K}[u]\Big|
  \le
  \|v\|_{L^2(\Omega)}\|\mathcal{K}[u]\|_{L^2(\Omega)}
  \le
  C_s\|u\|_{L^2(\Omega)}\|v\|_{L^2(\Omega)}.
\end{equation*}
The final inequality in the theorem follows from $\|w\|_{L^2(\Omega)}\le \mu^{-1/2}\|w\|_{\mathcal{U}}$ for $w=u,v$.
\end{proof}

\end{document}
